% !TEX root = ../notes.tex

\documentclass[../notes.tex]{subfiles}

\begin{document}

\section{March 4}
Here we go.

\subsection{The Master Theorem}
Let's try to upgrade our proof of \Cref{thm:slope}. Recall $\op{SynSp}\coloneqq\mathrm{Mod}(\op{fil}^*\mathbb S_{(p)})$, and there is a functor $\nu\colon\mathrm{Spectra}\to\mathrm{SynSp}$ given by
\[\nu(X)\coloneqq\lim\tau_{\ge2*}(X\otimes\mathrm{BP}^{\otimes(\bullet+1)}).\]
\begin{remark}
	If $X$ is bounded below, then $\nu(X)$ is $\tau$-complete, and $\nu X\left[\tau^{-1}\right]=X$. The moral is that we have been given a ``motivic lift'' from $X$ as a spectrum to a synthetic spectrum.
\end{remark}
\begin{remark}
	It turns out that $\nu$ commutes with filtered colimits of bounded below spectra. However, it does not preserve co\-fibers: the problem is that a map $f\colon X\to Y$ produces a sequence
	\[\tau_{2*}X\to\tau_{2*}Y\to\tau_{2*}\op{cofiber}(X\to Y)\]
	which is not expected to be a cofiber sequence. In particular, one has an issue with the numerology of the cutoff, so the difference from the above being a cofiber sequence is just given by an Eilenberg--MacLane space. As a result, one finds that the cofiber of the sequence
	\[\op{cofiber}(\nu f)\to\nu\op{cofiber}(f)\]
	is $\tau$-torsion.
\end{remark}
\begin{theorem} \label{thm:master-thm}
	There is a curve $C$ in the plane with the following property: for any $X$ with $\pi_iX=0$ for all $i<0$, then $\pi_{**}(\nu X)$ is $\tau$-power-torsion above the curve. Additionally, all lines of constant slope are eventually above $C$.
\end{theorem}
\begin{proof}
	Any such $X$ is a filtered colimit of spheres, so this follows from properties of $\nu$ and \Cref{thm:slope}. In particular, even though $\nu$ does not commute with cofibers, it does commute with them up to some $\tau$-torsion, which we notably do not care about for the application.
\end{proof}
Professor Hahn called \Cref{thm:master-thm} the ``master theorem,'' from which we will derive many other results of Devinatz--Hopkins--Smith.
\begin{remark} \label{rem:x-axis-adams-novikov}
	Let $I$ be the fiber of $\mathbb S_{(p)}\to\mathrm{BP}$. If $X$ is bounded below, then there is a spectral sequence associated to
	\[\cdots\to I^{\otimes2}\otimes X\to I\otimes X\to X\to X\to X\to\cdots.\]
	(The spectral sequence converges because $I$ is concentrated in positive degrees, so higher tensor-powers of $I$ only live in high degrees.) It turns out that the $E_2$-page of this spectral sequence recovers the one from $\pi_{**}\op{fil}^*\mathbb S/\tau$. The moral is that this gives an interpretation of the $x$-axis of our Adams--Novikov spectral sequence: on the $E_\infty$-page, the $x$-axis of $\pi_{**}\nu X$ is the image of $\pi_*X\to\mathrm{BP}_*X$!
\end{remark}
This interpretation of the Adams--Novikov spectral sequence has the following application.
\begin{corollary} \label{cor:connective-to-nilpotent-thing}
	Let $R$ be a $p$-local connective $\mathbb E_1$-ring (meaning that $R$ has no negative homotopy groups). Then every class in
	\[\ker(\pi_*R\to\mathrm{BP}_*R)\]
	is nilpotent.
\end{corollary}
\begin{proof}
	Let's look at the $E_\infty$-page of $\nu(R)$. Then by \Cref{rem:x-axis-adams-novikov}, every class in the kernel of the given map is a class strictly above the $x$-axis. But all such classes are nilpotent by \Cref{thm:master-thm}.
\end{proof}
\begin{remark}
	This is a first indication that \Cref{thm:slope} is going to be powerful: we have managed to prove something about rather general rings.
\end{remark}
We can actually upgrade \Cref{cor:connective-to-nilpotent-thing} to work with arbitrary $\mathbb E_1$-rings. This requires a trick.
\begin{definition}[smash nilpotent]
	Fix a $p$-local spectrum $X$. A class $\alpha\in\pi_iX$ is \textit{smash nilpotent} if and only if $\alpha^{\otimes N}\colon S^{iN}\to X^{\otimes N}$ is nullhomotopic for some $N$.
\end{definition}
\begin{proposition}
	Fix a spectrum $X$ and a class $\alpha\in\pi_iX$. Then $\alpha$ is smash nilpotent if $\mathrm{MU}\otimes\alpha$ is nullhomotopic.
\end{proposition}
\begin{proof}
	Note that any spectrum is filtered colimit of perfect spectra, so we can present $X$ as a filtered colimit of $\mathbb S_{(p)}$-modules, which we denote $X=\colim_iX_i$. Because $S^i$ is compact, we see that $S^i$ factors through one of those $X_i$. We may even factor the nullhomotopy of $\mathrm{MU}\otimes\alpha$ through one of the $X_i$s. Thus, we may reduce to the case where $X$ is a perfect $\mathbb S_{(p)}$-module. Because $X$ is now finitely constructed from $\mathbb S_{(p)}$s, it is bounded below, so $\Sigma^dX$ is connective for sufficiently large $d$. Furthermore, $\Sigma^d\alpha\in\pi_{i+d}\Sigma^dX$ will be nilpotent if and only if $\alpha$ is.

	We are thus basically reduced to the case where $X$ is connective. Note that we can form an $\mathbb E_1$-ring
	\[\mathbb S\oplus\Sigma^dX\oplus\left(\Sigma^dX\right)^{\otimes2}\oplus\cdots.\]
	(This is the free $\mathbb E_1$-ring on $\Sigma^dX$.) Being smash nilpotent is simply the claim that the image of $\Sigma^d\alpha$ in this $\mathbb E_1$-ring is nilpotent. But our class is now nilpotent in this ring by \Cref{cor:connective-to-nilpotent-thing}.
\end{proof}
\begin{corollary}
	Let $R$ be a $p$-local $\mathbb E_1$-ring. Then every class in
	\[\ker(\pi_*R\to\mathrm{BP}_*R)\]
	is nilpotent.
\end{corollary}
\begin{proof}
	Any element in the kernel of $S^i\to R$ is smash nilpotent. Taking a power is then nullhomotopic, so the original element must have been nilpotent.
\end{proof}
\begin{example} \label{ex:test-ring-zero-bp}
	Let $R$ be a $p$-local $\mathbb E_1$-ring. Then $R$ vanishes if and only if $1$ is nilpotent, which is then true if and only if $1$ is in the kernel of the map $\pi_*R\to\mathrm{BP}_*R$. Unwinding, we see that $R$ vanishes if and only if $\mathrm{BP}_*R$ vanishes.
\end{example}

\subsection{The Spectrum \texorpdfstring{$\mathrm{MUP}$}{ MUP}}
For our next application, we want the following construction.
\begin{definition}
	Let $\mathrm{MUP}$ be the Thom spectrum of the canonical map $\mathrm{BU}\times\ZZ\to\op{Pic}\mathbb S$.
\end{definition}
\begin{remark}
	It turns out that $\pi_*\mathrm{MUP}=\ZZ[\{x_i\}][u,1/u]$, where $u$ lives in degree $2$, and $x_i$ lives in degree $2i$.
\end{remark}
Recall that we can construct $\op{Pic}\mathrm{MUP}$, which is the subcategory of $\otimes$-invertible modules, equipped with only automorphisms as morphisms. There is a full subcategory $\op{BGL}_1\mathrm{MUP}\subseteq\op{Pic}\mathrm{MUP}$ above $\mathrm{MUP}$. For example,
\[\pi_i\op{BGL}_1\mathrm{MUP}\subseteq\op{Pic}\mathrm{MUP}=\begin{cases}
	0 & \text{if }i=0, \\
	\pi_0\mathrm{MUP}^\times & \text{if }i=1, \\
	\pi_{i-1}\mathrm{MUP} & \text{if }i\ge2.
\end{cases}\]
The case of $i=0$ has no content, the case of $i\ge2$ follows because the higher homotopies are all invertible in any $\infty$-category anyway. The case of $i=1$ holds because we are hunting for invertible morphisms.
\begin{example}
	A map $f\colon S^3\to\op{BGL}_1\mathrm{MUP}$ is given by an element in $\pi_2\mathrm{MUP}$; for concreteness, one could take $u\in\pi_*\mathrm{MUP}$. It turns out that $\op{Thom}(f)$ is some colimit of $S^3$, and by imagining $S^3$ as gluing together the two caps, we find that the Thom spectrum is
	\[\op{Thom}(f)=\op{cofiber}\left(f\colon\Sigma^2\mathrm{MUP}\to\mathrm{MUP}\right).\]
	Here is something amusing: $S^3$ can be realized as $\op{SU}(2)$, so one finds that $S^3=\Omega\mathbb{HP}^\infty$. Then $f$ lifts to a map $S^4=\Sigma S^3\to\mathrm B^2\op{GL}_1\mathrm{MUP}$. One may hope to extend this map from $S^4$ all the way to $\mathrm{HP}^\infty\to\mathrm B^2\op{GL}_1(\mathrm{MUP})$, which is possible because the only obstructions to gluing would live in the odd homotopy groups of the target, which are the even homotopy groups of $\mathrm{MUP}$, which all vanish. Thus, $f$ arises from loops of some map $\mathbb{HP}^\infty\to\mathrm B^2\op{GL}_1(\mathrm{MUP})$, so the cofiber $\op{Thom}(f)=\op{MUP}/f$ is an $\mathbb E_1$-algebra! Because the lifting is not unique, the $\mathbb E_1$-structure is not unique.
\end{example}
\begin{remark} \label{rem:quotient-by-pi2}
	In fact, one can iterate the previous example for any sequence of elements $f_1,f_2,\ldots\in\pi_2\mathrm{MUP}$ to show that $\mathrm{MUP}/(f_1,f_2,\ldots)$ is an $\mathbb E_1$-algebra. For example, the quotient
	\[R\coloneqq\frac{\mathrm{MUP}}{\left(2u,x_1-u,x_2u^{-1},x_3u^{-2},\ldots\right)}\]
	admits the structure of an $\mathbb E_1$-algebra and has homotopy groups $\pi_*R=\FF_2[x_1,1/x_1]$, so it is a skew field.
\end{remark}
Thus, we have produced a rather interesting example of a skew field. We will shortly classify all of them.

\end{document}